
\begin{question}
	\questiontext{Triadic Closure and Tie Stability in Social Networks}
	\begin{subquestion}{Analyze the structural configuration of node A with respect to the Strong Triadic Closure (STC) principle. Discuss under what changes to the network structure or to the strength of existing ties this configuration could become consistent with the STC principle.}
		\answer{
			To satisfy the STC principle since  A has strong ties to B and C, then there should be an edge (either a strong or weak tie) between B and C. since the question has stated that AD is a local bridge than by definition since A should satisfy STC property the tie is a weak one so no more modifications are needed but if the edge AD wasn't a local bridge and it was strong tie based on STC principle there should have been an edge between B and D (either a strong or weak tie) and also an edge Between C and D  (either a strong or weak tie)
		}
	\end{subquestion}
	\begin{subquestion}{Assume that the network satisfies the Strong Triadic Closure principle. Using a logical argument and the definition of a local bridge, show that the edge (A, D) cannot be a strong tie.
		}
		\answer{
		In this case we will use a proof by contradiction assume for the sake of contradiction that the edge $A \to D$ is a strong tie so by definition of the Strong Triadic Closure principle and edge of any strength should exist between nodes $D$ and $C$ and also an edge of any strength should exist between nodes $D$ and $B$ which contradicts the fact that the edge $A \to D$ is a local bridge therefore, our assumption is false and hence the edge  $A \to D$ is a weak tie
		}
	\end{subquestion}
\end{question}